\begin{resumo}[RESUMO]
\begin{SingleSpace}
\noindent Este documento apresenta o plano de crescimento e escalabilidade da startup PetCuida, plataforma digital de prevenção e acesso a cuidados para animais de companhia. O diagnóstico situa o negócio na fase de validação inicial, ainda sem receita nem usuários pagantes, e identifica como principais desafios a disposição a pagar não medida, a dependência de convênios para o resultado positivo, o custo do fundo de liquidação de créditos solidários, que supera a contribuição unitária, o custo de aquisição não mensurado e o arranque simultâneo dos dois lados da plataforma. Propõem-se duas estratégias de crescimento: a densidade local por meio de parcerias com clínicas e organizações, expandindo-se por praças somente após gatilhos mensuráveis, e o aprofundamento de receita com novos serviços condicionados a evidências. O primeiro ciclo da estratégia principal é um piloto territorial assistido de oito semanas, com definições operacionais de ativação e retenção e critérios de decisão. As ações de escalabilidade abrangem automação, uso de tecnologia, modelo replicável por praça e receita recorrente. Uma simulação baseada na economia unitária da Atividade IX mostra que a escala, por si só, é insuficiente: a contribuição de R\$~0,55 por usuário ativo por mês só se torna robusta com alavancas combinadas. O plano termina com quadro de ação, metas, indicadores e uma análise explícita de viabilidade e coerência.

\vspace{\onelineskip}
\noindent\textbf{Palavras-chave}: startup; estratégia de crescimento; escalabilidade; plataforma de dois lados; economia unitária.
\end{SingleSpace}
\end{resumo}
